\documentclass[10pt,a4paper]{amsart}
\usepackage[margin=19mm]{geometry}
\usepackage{amsmath,amssymb,mathtools}
\usepackage[scr=rsfs,cal=euler]{mathalfa}
\usepackage{xfrac}
\usepackage[unicode,hidelinks,bookmarks=false]{hyperref}
\newcommand{\ie}{i.e.\ }
\newcommand{\cf}{cf.\ }
\newcommand{\resp}{respectively\ }
\newcommand{\Subsection}{\subsection}
\providecommand{\nobreakdash}{\nobreak\hbox{-}}
\setlength{\emergencystretch}{2em}
\multlinegap0pt
\usepackage{amssymb}
\usepackage{enumerate}
\usepackage[shortlabels]{enumitem}
\usepackage{tikz}
\usepackage{float}
\usepackage{mathtools}
\usepackage{dsfont}
\usepackage[normalem]{ulem}
\newtheorem{lemma}{Lemma}
\newcommand{\ZZ}{{\mathbb {Z}}}
\newcommand{\res}{{\rm res}}
\title{Monoidal categorification of generalized cluster algebras and conjectures of Fraser and Gleitz}
\author{Xiao-Juan An, Jian-Rong Li, and Yan-Feng Luo$^\dag$}
\date{Source: arXiv:2606.16361v2 (CC BY 4.0)}
\begin{document}
\maketitle

\noindent\emph{Source and licence.} Monoidal categorification of generalized cluster algebras and conjectures of Fraser and Gleitz. Version arXiv:2606.16361v2 (CC BY 4.0), distributed by arXiv under the Creative Commons Attribution 4.0 International licence, http://creativecommons.org/licenses/by/4.0/, as recorded in the arXiv licence metadata for this exact version and shown on the abstract page https://arxiv.org/abs/2606.16361v2. Full licence notice: \texttt{licenses/arxiv-2606.16361-CC-BY-4.0.txt}.\par\smallskip

\noindent\textit{Editorial scope.} Counted unit: the article's lemma special (source0.tex lines 1105--1111). The article's complete original proof of it is delivered with it (source0.tex lines 1113--1169, one \texttt{proof} environment of 57 lines). No mathematics is written, repaired or re-derived: the statement and the proof are the article's own lines, in the article's own notation, and no retained line is substituted or re-flowed.  Retained uncounted as the source-local context the proof needs: lines 463--465; lines 1087--1089.  Nothing was omitted from any retained span, so the package carries no omission record.  No retained line contains a citation, so the wrapper carries no bibliography and no bibliography entry.  No retained line contains a figure environment, an image-inclusion command or drawing code, so no image is delivered and the package declares no asset dependency.  Editorial formatting only, in addition: an AMS \texttt{amsart} preamble that restates the article's own declarations and macros used by the retained lines, namely the environments \texttt{lemma} on separate counters, together with the standard packages those lines need.  The retained environments print in this document as Lemma 1 in order of appearance, whereas the article prints the same items with the numbering of its own full document.  The complete native source of the article as distributed in the arXiv e-print for this exact version is bundled unchanged as \texttt{sources/202931/source0.tex}.
\begin{align}\label{Non-overlapping paths}
\overline{\mathscr{P}}_{(i_{t},k_{t})_{1\leq t\leq z}}=\{(p_{1},\ldots,p_{z}): p_{t}\in \mathscr{P}_{i_{t},k_{t}}, 1\leq t\leq z, (p_{1},\ldots,p_{z})\text{ is } \text {non-overlapping}\}.
\end{align}

\begin{align}\label{R}
R_{(i,k,s)(j,v,t)}=\sum_{(p_{1},\ldots,p_{s+t}) \in \overline{\mathscr{P}}_{((i,k+2r)_{0 \leq r \leq s-1}(j,v+2r)_{0 \leq r \leq t-1})}} \prod_{z=1}^{s+t}\mathfrak{m}(p_{z}),
\end{align}

\begin{lemma}\label{Lem: minimal affine is special}
For $U_{\varepsilon}^{\res}({L\mathfrak{sl}_{3}})$, let $\varepsilon^{2 \ell}=1$ with $\ell\in \ZZ_{\geq 2}$. Let $\{i,j\}=\{1,2\}$ and let $(i,k), (j,v)\in\mathcal{X}$ with $v=k+2 \ell-1$. Then, for $1\leq t \leq \ell-1$, 
\begin{align*}
\chi_\varepsilon(L(M^{j,v,t}_{i,k,\ell-1}))=\chi_q(L(M^{j,v,t}_{i,k,\ell-1}))|_{q=\varepsilon}.
\end{align*}
Moreover, for $t=\ell-1$, the minimal affinizations $L(M^{j,v,t}_{i,k,\ell-1})$ are special.
\end{lemma}

\begin{proof}
We prove the lemma in the case $i=1$, $j=2$, and $k=0$; the proofs for the other cases follow a similar approach. For $t\leq \ell-1$, we need to prove the identity 
\[
\chi_\varepsilon(L(M^{2, 2\ell-1,t}_{1,0,\ell-1}))=\chi_q(L(M^{2,2\ell-1,t}_{1,0,\ell-1}))|_{q=\varepsilon},
\]
that is,
{\small
\begin{align*}
\chi_\varepsilon(L(Y_{1,0}Y_{1,2}\cdots Y_{1,2\ell-4}Y_{2,2\ell-1}Y_{2,2\ell+1}\cdots Y_{2,2 \ell+2t-3}))=\chi_q(L(Y_{1,0}Y_{1,2}\cdots Y_{1,2\ell-4}Y_{2,2\ell-1}Y_{2,2\ell+1}\cdots Y_{2,2\ell+2t-3}))|_{q=\varepsilon}.
\end{align*}}
To establish the above identity, it suffices to show that $R_{(1,0,\ell-1)(2,2\ell-1,t)}$ contains no dominant monomial $m$ except the highest weight monomial such that the $\varepsilon$-character of $L(m)$ is contained in $R_{(1,0,\ell-1)(2,2\ell-1,t)}$. Monomials in $R_{(1,0,\ell-1)(2,2\ell-1,t)}$ are of the form $\mathfrak{m}(p_1)\cdots \mathfrak{m}(p_{t+\ell-1})$, where  $p_i \in \mathscr{P}_{1,2i-2}$ for $i\in[1, \ell)$, and $p_i \in \mathscr{P}_{2,2i-1}$ for $i\in[\ell, t+\ell-1]$, see Formula (\ref{R}). The proof is divided into the following three cases.

{\bf Case} 1. Suppose that $p_a\in \mathscr{P}_{1,2a-2}$ with $\mathfrak{m}(p_a)=Y_{1,2a-2}$ for some $a\in[1, \ell)$. To obtain a dominant monomial, a path $p_b\in \mathscr{P}_{2,2b-1}$, for some $b\in[\ell,\ell+t-1]$, must either consist entirely of upper corners, or have $(1,2a-2+2 \ell)$ as its unique lower corner. If $p_b$ consists entirely of upper corners, then $\mathfrak{m}(p_b)=Y_{2,2b-1}$, where $b\in[\ell,\ell+t-1]$. In this case, we obtain the dominant monomial $Y_{1,0}Y_{1,2}\cdots Y_{1,2\ell-4}Y_{2,2\ell-1}Y_{2,2\ell+1}\cdots Y_{2,2\ell+2t-3}$, which is the highest weight monomial. If $p_b$ (with $b\in[\ell,\ell+t-1]$) has $(1,2a-2+2\ell)$ as its unique lower corner, then $a\geq 2$ and $\mathfrak{m}(p_b)=Y^{-1}_{1,2a-2+2 \ell}=Y^{-1}_{1,2b+2}$, where $p_b\in\mathscr{P}_{2,2b-1}$. By Formula (\ref{Non-overlapping paths}), we have $\mathfrak{m}(p_u)=Y^{-1}_{1,2u+2}$, where $p_u\in\mathscr{P}_{2,2u-1}$, $u\in [b, \ell+t-1]$.

In other words, in order to obtain a dominant monomial, for $a\in[1,\ell)$, we take $\mathfrak{m}(p_a)=Y_{1,2a-2}$. For $z\in[\ell,b)$, we take $\mathfrak{m}(p_z)=Y_{2,2z-1}$. For $u\in[b, \ell+t-1]$, we take $\mathfrak{m}(p_u)=Y^{-1}_{1, 2u+2}$. Thus, we obtain the monomial 
\begin{align}\label{monomial}
Y_{1,0}\cdots Y_{1,2 \ell-4}Y_{2,2\ell-1}\cdots Y_{2,2b-3} Y^{-1}_{1,2b+2}\cdots Y^{-1}_{1,2 \ell+2t}.
\end{align}

If $t=\ell-1$, then, since $a\geq 2$, the factor $\mathfrak{m}(p_{\ell+t-1})=Y^{-1}_{1,2 \ell+2t}=Y^{-1}_{1,4 \ell-2}$ cannot be canceled by any other factor appearing in the product. Therefore, the monomial (\ref{monomial}) is not dominant.

If $t<\ell-1$, then the monomial (\ref{monomial}) simplifies to 
\[
Y_{1,0}\cdots Y_{1,2a-4}Y_{1,2a+2t+2}\cdots Y_{1,2\ell-4}Y_{2,2 \ell-1}\cdots Y_{2,2b-3}, \quad 
\]
which is a dominant monomial. Its lowest weight monomial is  
\[
Y^{-1}_{2,3} Y^{-1}_{2,5}\cdots Y^{-1}_{2,2 \ell-1} Y^{-1}_{1,2 \ell+2}\cdots Y^{-1}_{1,2b} \cdot Y_{2,2b-1}\cdots Y_{2,2 \ell+2t-3},
\]
which does not appear in $R_{(1,0,\ell-1)(2,2\ell-1,t)}$, since $Y^{-1}_{1,2b}$ and $Y_{2,2b-1}$ are overlapping. Thus, there is no dominant monomial $m$ except the highest weight monomial such that the $\varepsilon$-character of $L(m)$ is contained in $R_{(1,0,\ell-1)(2,2\ell-1,t)}$.


{\bf Case} 2. Suppose that $p_a\in \mathscr{P}_{1,2a-2}$ with $\mathfrak{m}(p_a)=Y^{-1}_{1,2a}Y_{2,2a-1}$ for $a\in[1, \ell)$. To obtain a dominant monomial, a path $p_b\in \mathscr{P}_{2,2b-1}$, for $b\in[\ell,t+\ell-1]$, must have $(1,2a+2+2 \ell)$ as its upper corner. In this case, we have 
\[
\mathfrak{m}(p_b)=Y_{1,2a+2 \ell}Y^{-1}_{2,2a+2 \ell+1}=Y_{1,2b}Y^{-1}_{2, 2b+1}.
\]
Hence
\[
\mathfrak{m}(p_a)\mathfrak{m}(p_b)=Y^{-1}_{1,2a}Y_{2,2a-1}Y_{1,2a+2 \ell}Y^{-1}_{2,2a+2 \ell+1}=Y_{2,2a-1}Y^{-1}_{2,2a+2 \ell+1}.
\]
To obtain a dominant monomial, let $\mathfrak{m}(p_{a+1})=Y^{-1}_{1,2a+2}Y_{2,2a+1}$. Therefore,
\[
\mathfrak{m}(p_a)\mathfrak{m}(p_b)\mathfrak{m}(p_{a+1})=Y_{2,2a-1}Y^{-1}_{2,2a+2 \ell+1}Y^{-1}_{1,2a+2}Y_{2,2a+1}=Y_{2,2a-1}Y^{-1}_{1,2a+2}.
\]
By Formula (\ref{Non-overlapping paths}), proceeding in this way, since $t\leq \ell-1$, we will always find that there exists a monomial, namely $Y^{-1}_{1,2 \ell-2}$ or $Y^{-1}_{2,2\ell+2t-1}$, which cannot be canceled. Thus, in this case, no dominant monomial appears in $R_{(1,0,\ell-1)(2,2\ell-1,t)}$.

{\bf Case} 3. Suppose that $p_a\in \mathscr{P}_{1,2a-2}$ with $\mathfrak{m}(p)=Y^{-1}_{2,2a+1}$ for some $a\in[1, \ell)$. To obtain a dominant monomial, the path $p_b\in \mathscr{P}_{2,2b-1}$, for some $b\in[\ell,\ell+t-1]$, must have $(2,2a+2\ell+1)$ as its upper corner. Then 
\[
\mathfrak{m}(p_b)=Y_{2,2a+2\ell+1}=Y_{2,2b-1}.
\]
Using the same method as in Case 1, we conclude that no dominant monomial appears in $R_{(1,0,\ell-1)(2,2\ell-1,t)}$ in this case.

In summary, we have
\begin{align*}
\chi_\varepsilon(L(M^{j,v,t}_{i,k,\ell-1}))=\chi_q(L(M^{j,v,t}_{i,k,\ell-1}))|_{q=\varepsilon}.
\end{align*}
\end{proof}  

\end{document}
